\subsection{Flat divisors and Hartogs extension}
\label{subsec:flat-divisor-complements}

All the nine-point blowups of \cref{prop:nine-point-anticanonical-curvature} carry backgrounds of order one, so the order of degeneration does not distinguish the torsion case, in which a multiple of $D$ moves in a basepoint-free pencil, from the non-torsion case, in which no multiple of $D$ moves. The difference is visible in the complex geometry near $D$. On the tube where the curvature form of \cref{prop:flat-divisor-rank-one-curvature} vanishes, the logarithm of the flat norm of the defining section is pluriharmonic. Its regular level sets are therefore Levi-flat, that is, real hypersurfaces foliated by complex hypersurfaces. These complex leaves are governed by the unitary holonomy of the normal bundle: finite holonomy gives compact fibers of a local holomorphic function, whereas infinite holonomy makes every leaf dense and forces every holomorphic function on the punctured tube to be constant. In both cases the complement of $D$ is not Stein.

In this subsection $S$ is a complex manifold of dimension $n\geq2$ and $D\subset S$ is a smooth compact connected divisor. Suppose that $L=\mathcal O_S(D)$ has a unitary flat metric $h$ on a neighborhood of $D$. Write $s_D$ for its defining section and $N=L|_D$ for the normal bundle, with the induced metric and flat connection, and let $\chi\colon\pi_1(D)\to\mathrm U(1)$ be the holonomy character of $N$, which describes parallel transport around closed loops. In parallel unit frames, $s_D$ is represented by holomorphic functions $W_j$ satisfying $W_j=t_{jk}W_k$ with locally constant $t_{jk}\in\mathrm U(1)$. The flat connection defines on each circle bundle $\{v\in N:|v|=r\}$ a foliation whose leaves are given locally, in parallel frames, by $\{w=\text{constant}\}$; we call it the flat horizontal foliation.

\Needspace{16\baselineskip}
\begin{lemma}\label{lem:flat-divisor-foliated-tube}
There are $R>0$, a neighborhood $U$ of $D$ and a smooth diffeomorphism
\[
 \Upsilon\colon\{v\in N:|v|<R\}\longrightarrow U
\]
that is the identity on the zero section and preserves the defining function: in each parallel normal frame, $W_j(\Upsilon(p,w))=w$. In particular, $|s_D|_h\circ\Upsilon=|w|$. For $0<r<R$, the set
\[
 M_r=\{x\in U:|s_D(x)|_h=r\}
\]
is a compact connected Levi-flat hypersurface, and $\Upsilon$ maps the leaves of the flat horizontal foliation of $\{|w|=r\}$ onto the complex leaves of $M_r$.
\end{lemma}
\begin{proof}
Choose a smooth normal exponential map $\Upsilon_0$ whose first normal derivative is complex linear; the derivative of $s_D$ identifies the normal bundle with $L|_D$, and we use this identification to fix that first derivative. Parallel transport along the normal rays identifies $\Upsilon_0^*L$ with the pullback of $N$. In compatible parallel frames one then has
\[
 W_j(\Upsilon_0(p,a))=a+O(|a|^2).
\]
For each $p$, this is a smooth map between real two-dimensional disks whose derivative at zero is the identity. The inverse function theorem with parameters, uniformly over the compact base $D$, gives a unique small $a=a(p,w)$ with $W_j(\Upsilon_0(p,a))=w$. It depends smoothly on $(p,w)$, and its first derivative in $w$ at $w=0$ is the identity. On an overlap of frames, $a$ and $w$ change by the same constant unitary factor, so uniqueness makes these corrections agree. After decreasing $R$, the map $(p,w)\mapsto(p,a(p,w))$ is a diffeomorphism of a tube onto its image, and its composition with $\Upsilon_0$ is the required map $\Upsilon$.

Off $D$, the function $F=\log|s_D|_h^2$ is pluriharmonic, and its radial derivative in the tube is nonzero. Hence every $M_r$ is a smooth hypersurface; it is compact and connected, being a circle bundle over $D$. Locally its complex tangent space is $\ker\mathrm dW_j$, so it is Levi-flat, with complex leaves the hypersurfaces $\{W_j=\text{constant}\}$. The identity $W_j\circ\Upsilon=w$ pulls this distribution back to $\{\mathrm dw=0\}$ in parallel frames, which is the flat horizontal foliation of the circle bundle. The map $\Upsilon$ is only a smooth tubular identification; its base projection need not be holomorphic.
\end{proof}

\begin{proposition}\label{prop:flat-divisor-finite-holonomy}
In this setting, $N$ is torsion as a holomorphic line bundle if and only if its unitary holonomy has finite image. If $N$ has exact order $m$, then a sufficiently small tube $U$ admits a proper surjective holomorphic map
\[
 \phi\colon U\longrightarrow\{z\in\mathbb C:|z|<R^m\}
\]
with connected fibers and $\operatorname{div}(\phi)=mD$. Its fibers over nonzero values are smooth compact complex hypersurfaces.
\end{proposition}
\begin{proof}
If $\chi^m=1$, a parallel nowhere-zero section trivializes $N^m$ holomorphically. Conversely, if a holomorphic section $e$ trivializes $N^m$, then $\log|e|^2$ is pluriharmonic on the compact connected manifold $D$, hence constant by the maximum principle. In parallel holomorphic frames the coefficient of $e$ is a holomorphic function of constant modulus, hence constant. Therefore $e$ is parallel and $\chi^m=1$. Consequently the least such $m$ is both the order of $N$ as a torsion line bundle and the order of the finite holonomy group.

By \cref{lem:flat-divisor-foliated-tube}, the tube $U$ retracts onto $D$. Hence the flat bundle $L^m$ on $U$ has trivial holonomy and a global parallel section $e_U$ of unit length. Define $\phi=s_D^m/e_U$. Then $\phi$ is holomorphic, $\operatorname{div}(\phi)=mD$ and $|\phi|=|s_D|_h^m$. The inverse image of a compact subset of the target lies in a smaller closed normal disk bundle and is closed there, which proves properness, and the tubular model shows that every value in the disk is attained.

Choose local parallel unit frames $e_j$ of $L$ with $e_j^m=e_U$, so that $\phi=W_j^m$. In the tubular model, the fiber over a nonzero value consists, over each point of $D$, of the $m$ roots of that value. Parallel transport acts on these roots through $\chi$, whose image has exact order $m$ and therefore acts transitively. Thus this fiber is a connected covering of $D$, as a real manifold. Since $\mathrm d\phi=mW_j^{m-1}\mathrm dW_j$, it is a smooth complex hypersurface, and it is compact because $\phi$ is proper. The fiber over zero has connected support $D$, which completes the proof that all fibers are connected.
\end{proof}

\begin{theorem}\label{thm:flat-divisor-holomorphic-rigidity}
\index{Hartogs extension!non-torsion flat divisor}
\index{Levi-flat hypersurface!flat divisor tube}
Suppose $N$ is not torsion. Then every complex leaf of $M_r$ is dense in $M_r$, and every holomorphic function on $U\setminus D$ is constant:
\[
 H^0(U\setminus D,\mathcal O_S)=\mathbb C.
\]
Suppose moreover that $S$ is connected. For every compact set $K\subset S\setminus D$ such that $(S\setminus D)\setminus K$ is connected, every holomorphic function on $(S\setminus D)\setminus K$ is constant, and in particular extends holomorphically to $S\setminus D$. Similarly, for every connected open neighborhood $V$ of $D$, every holomorphic function on $V\setminus D$ is constant and extends to $S\setminus D$.
\end{theorem}
\begin{proof}
By \cref{prop:flat-divisor-finite-holonomy}, $\chi(\pi_1(D))$ is an infinite subgroup of $\mathrm U(1)$. Such a subgroup is dense, since a proper closed subgroup of the circle is finite, as one sees by lifting it to a closed subgroup of $\mathbb R$ containing $2\pi\mathbb Z$. The circle bundle with its flat horizontal foliation is the suspension of $\chi$, that is, the quotient of $\widetilde D\times S^1$ by the diagonal action of $\pi_1(D)$, acting on $S^1$ through $\chi$. A leaf projects onto $D$, and its intersection with each fiber contains the holonomy orbit of any one of its points, which is dense in the fiber. In a local trivialization, horizontal transport maps the orbits in nearby fibers to each other by a fixed rotation, so the leaf is dense in the whole circle bundle. By \cref{lem:flat-divisor-foliated-tube}, every complex leaf is dense in $M_r$.

Let $f$ be holomorphic on $U\setminus D$. Its modulus attains a maximum on the compact set $M_r$. On the complex leaf through a maximum point, $|f|$ is bounded by this maximum and attains it, so the maximum modulus principle makes $f$ constant on the connected leaf. By density and continuity, $f$ equals the same constant $c$ on $M_r$. Near any point of $M_r$, choose holomorphic coordinates $(z,W_j)$ with $z\in\mathbb C^{n-1}$. For each fixed $z$, the holomorphic function $W_j\mapsto f(z,W_j)-c$ vanishes on an arc of the circle $|W_j|=r$, so it vanishes on the coordinate neighborhood by the one-variable identity theorem. Analytic continuation on the connected punctured tube gives $f\equiv c$.

A compact set $K\subset S\setminus D$ is disjoint from a sufficiently small tube around the compact divisor $D$. A holomorphic function on $(S\setminus D)\setminus K$ is therefore constant on this punctured tube, and by the identity theorem on the connected set $(S\setminus D)\setminus K$ it is constant everywhere; the constant function is the required extension. Similarly, every open neighborhood $V$ of $D$ contains a full tube, and $V\setminus D$ is connected, because paths in the connected manifold $V$ can be perturbed off the real codimension-two submanifold $D$. The identity theorem gives the assertion for $V$. The same path argument shows that $S\setminus D$ is connected.
\end{proof}

Recall that a complex manifold $Y$ is \emph{holomorphically separable} if its holomorphic functions separate points, and \emph{holomorphically convex} if the holomorphic hull
\[
 \widehat K_Y=\{y\in Y:|f(y)|\leq\max_K|f|\text{ for every }f\in\mathcal O(Y)\}
\]
of every compact set $K\subset Y$ is compact. A Stein manifold has both properties.

\begin{corollary}\label{cor:flat-divisor-nonstein}
\index{Stein manifold!flat divisor obstruction}
For any holonomy of $N$, the punctured tube $U\setminus D$ admits no smooth strictly plurisubharmonic function, and $S\setminus D$ is not Stein. In the non-torsion nine-point blowups of \cref{prop:nine-point-anticanonical-curvature}, every holomorphic function on the complement of $D$ is constant, and the Hartogs extension of \cref{thm:flat-divisor-holomorphic-rigidity} across compact holes holds.
\end{corollary}
\begin{proof}
The restriction of a smooth strictly plurisubharmonic function on $U\setminus D$ to a complex leaf of $M_r$ is strictly plurisubharmonic. At a maximum point on the compact set $M_r$, this restriction has a local maximum on the leaf, which has positive dimension because $n\geq2$; this contradicts the maximum principle.

In the finite-holonomy case, each fiber of the map of \cref{prop:flat-divisor-finite-holonomy} over a nonzero value is a connected compact complex hypersurface of positive dimension. By the maximum modulus principle, every holomorphic function on $S\setminus D$ is constant on it, so two distinct points of the fiber cannot be separated. In the infinite-holonomy case, \cref{thm:flat-divisor-holomorphic-rigidity} makes every holomorphic function constant on the punctured tube, which again contradicts holomorphic separability. In both cases $S\setminus D$ is not Stein. Finally, \cref{thm:flat-anticanonical-neighborhood} gives the flat neighborhood for the nine-point blowups, whose normal bundle is non-torsion by assumption, and \cref{thm:flat-divisor-holomorphic-rigidity}, also with $K=\varnothing$, gives the last assertion.
\end{proof}

The nine-point blowups show that a background with a well-defined order of degeneration can have very different holomorphic geometry near its null locus; \cref{app:cotangent-torsors} gives a second instance of this contrast, for K\"ahler classes on a ruled surface over an elliptic curve.
